\chapter[Introduction]{Introduction}%[Chapter name as it will appear in the ToC]{chapter name as it will apear as the title to the chapter} Probably always want them to be the same 
\label{ch:introduction}

\section{Thesis to-do}

Write your thesis! Keep going to see some examples of thesis content: some writing with references, equations, quantities with SI units, overlong and sideways tables, among other things. Note that, like for previous chapter-level entries, there is a blank page after this to ensure that your chapter starts on an odd-numbered page.

\section{Creating a diff file for corrections}
This is a bit of a pain. I'd recommend:
\begin{enumerate}
    \item Using one of the suggested implementations of Overleaf \textcolor{blue}{\hyperlink{https://www.overleaf.com/learn/latex/Articles/How_to_use_latexdiff_on_Overleaf}{here}} --- I used the first one
    \item For both the new and old versions, copying your entire thesis into one .tex file, replacing the \texttt{input} and \texttt{include} commands
    \item Removing all references to the font in the preambles and changing the compiler to \texttt{pdflatex}
    \item Making sure that there are no instances of three or more blank new lines in the \texttt{.tex} files
    \item Having a folder for images that have stayed the same, and a folder for images that have changed, and making sure that new images point to the new folder.
    \item Debugging any errors that arise using the log and the \texttt{main-d.tex} file that is created (for the first solution) which is a combined \texttt{main} using both your files.
\end{enumerate}
